%---------------------------------------------------------------------
\chapter{Limitations, Ethics and Safety}
%---------------------------------------------------------------------

This chapter states the limitations of the data, the method and the evaluation, lists the ways in
which the agent can fail together with the guardrail that addresses each, and discusses the ethical
considerations of a system whose output affects producers and buyers.

\section{Limitations}

\paragraph{Data.} The producer-level evidence, the weaver registry and the batch records, is simulated.
Real individual records are not published, so the weaver and batch checks demonstrate the mechanism
but say nothing about the availability or quality of real records. The GI description passages for the
three original products were written by the team as summaries and their page references are
indicative; they were not extracted from the scanned GI application documents. For the twelve products
taken from the government GI catalogue, the descriptions are short and the catalogue does not give GI
application numbers. The registry stores one current state per product, so historical state names, such
as Andhra Pradesh for products now in Telangana, are not recognised.

\paragraph{Method.} All factual checks are deterministic string comparisons. As shown in
Table~\ref{tab:probes}, spelling variants of districts, state abbreviations and aliases that are not in
the alias table lead to wrong \emph{Conflicting} verdicts. Retrieval uses BM25 over 21 short passages,
which is adequate for this corpus but would not scale to the full text of GI applications without dense
retrieval and chunking. The language-model path for extraction and explanation is implemented but was
not evaluated, and the OCR path for label images was not installed or tested. The durable-workflow
support stores every state but does not resume interrupted runs automatically.

\paragraph{Evaluation.} The 45-case golden set and the 15 demonstration cases were constructed by the
team from the same rules and records that the agent uses, the labels were assigned by construction
rather than by independent annotators, and one rule was changed after the first evaluation run. The
100\% accuracy therefore shows that the agent behaves as designed on controlled cases; it is not an
estimate of accuracy on real marketplace listings. Citation faithfulness was not measured. The golden
set covers three of the fifteen products. The single-lookup baseline is weak by design.

\section{Failure Modes and Guardrails}

Table~\ref{tab:failures} pairs each identified failure mode with its mitigation.

\begin{table}[H]
\centering
\caption{Failure modes and mitigations}
\label{tab:failures}
\small
\begin{tabularx}{\linewidth}{@{}>{\raggedright\arraybackslash}X >{\raggedright\arraybackslash}X@{}}
\toprule
\textbf{Failure mode} & \textbf{Mitigation} \\ \midrule
Verdict issued without evidence & Citation gate: a decided verdict needs at least one source ID and a grade $\geq 0.6$, otherwise it is downgraded to \emph{Insufficient} and escalated \\
Weak or off-product retrieval & Grader detects unresolved product or weak search; planner re-queries with canonical name and product filter; maximum three iterations \\
Unbounded looping or tool calls & Iteration limit of 3 and a graph recursion limit of 40; observed maximum 2 iterations and 11 tool calls \\
Missing weaver or batch details & Checks marked \emph{unknown}, never \emph{pass}; case escalated to review with the missing items listed \\
Typing error in an identifier & Unknown identifier gives \emph{unknown}, not \emph{fail}; batch mismatch is not asserted against an unverified weaver \\
Spelling variants causing false conflicts (observed, Table~\ref{tab:probes}) & Not yet mitigated in the agent; planned fuzzy matching with an \emph{unknown} band, and abbreviation and former-name tables (Section~10.2) \\
Product outside the covered GIs & Interface reports ``not covered'' before the agent runs and offers referral to a reviewer \\
Model returns malformed or invented fields & JSON mode, Pydantic validation and deterministic fall-back parser; the model never performs factual checks \\
Instructions embedded in label text (prompt injection) & Label text is only parsed into fields; tools are fixed by the graph, not chosen by the model; the explanation prompt receives only check statuses, source IDs and tool notes \\
Reviewer error or bias & Reviewer must record a note; the decision is stored with the automated verdict, so both remain auditable \\
\bottomrule
\end{tabularx}
\end{table}

\section{Ethical and Responsible-Use Considerations}

\paragraph{Who could be harmed.} A wrong \emph{Conflicting} verdict on a genuine product could damage the
reputation and sales of a weaver who is legitimately entitled to the GI, and Chapter~7 showed that
spelling variants can produce exactly this error. A wrong \emph{Supported} verdict would mislead a buyer
and undermine the value of the GI for all producers in the region. The design therefore treats a false
\emph{Supported} verdict as the most serious error, and none occurred in the evaluation; the observed
false \emph{Conflicting} verdicts on realistic spelling variants are the most important weakness to
correct before any real use.

\paragraph{Advisory status.} Every verdict is labelled in the interface as ``advisory result based on
available records --- not a legal certification''. Authority over GI registration and authorised use
rests with the GI Registry under the GI Act~\cite{giact1999}, and the system does not claim to replace
it. Every verdict shows the checks performed and the sources used, so that a user can see why it was
reached, and any user can send a decided case to a human reviewer.

\paragraph{Human oversight.} Cases that the agent cannot decide are never resolved automatically. They
are routed to a reviewer, whose decision and note are stored separately from the automated verdict and
shown to the requester. This keeps a person accountable for every decision that is not fully supported
by evidence.

\paragraph{Privacy.} Weaver records can contain personal information. In this project, all weaver
records are simulated, and the names are randomly combined from lists of common first names and
surnames; they do not refer to real people. A real deployment would require consent from weavers and
cooperatives, access control on the weaver registry and the review dashboard, and a retention policy
for case data, which are not implemented here; the current Admin page has no authentication.

\paragraph{Fairness.} Coverage is uneven: fifteen GIs are supported, the alias table reflects the names
known to the team, and products from regions with less documentation are more likely to be marked
\emph{Insufficient} or ``not covered''. This should be reduced by extending coverage with producer
associations rather than by relaxing the checks.
